\documentclass[oneside,openany,phd]{TMU-Thesis}

% مشخصات پایان‌نامه را در فایلهای faTitle و enTitle وارد نمایید.
% در اینجا برای یکپارچگی، همینجا وارد می‌کنیم

% مشخصات فارسی
\university{تربیت مدرس}
\faculty{مهندسی برق و کامپیوتر}
\department{گروه هوش مصنوعی و رباتیکز}
\subject{مهندسی کامپیوتر}
\field{نرم‌افزار}
\credits{---} % در رساله دکتری معمولا به صورت واحدی ذکر نمی‌شود یا فیلد خالی است
\title{یادگیری تقویتی عمیق با شکل‌دهی تصادفی پاداش برای مدیریت عدم قطعیت در سیستم‌های خودوفقی}
\firstsupervisor{دکتر ...}
% \secondsupervisor{استاد راهنمای دوم}
\firstadvisor{دکتر ...}
% \secondadvisor{استاد مشاور دوم}
\name{علیرضا}
\surname{کاویانی‌فر}
\studentID{...}
\thesisdate{بهمن ۱۴۰۴}

% مشخصات انگلیسی
\latinuniversity{Tarbiat Modares University}
\latinfaculty{Faculty of Electrical and Computer Engineering}
\latinsubject{Computer Engineering}
\latinfield{Software}
\latintitle{Randomized Reward-Shaped Deep Reinforcement Learning for Managing Uncertainty in Self-Adaptive Systems}
\firstlatinsupervisor{Dr. ...}
\firstlatinadvisor{Dr. ...}
\latinname{Alireza}
\latinsurname{Kavianifar}
\latinthesisdate{February 2026}
\latinkeywords{Self-Adaptive Systems, Deep Reinforcement Learning, Reward Shaping, Uncertainty}

% فراخوانی بسته‌ها و دستورات استاندارد
% در این فایل، دستورها و تنظیمات مورد نیاز، آورده شده است.
%-------------------------------------------------------------------------------------------------------------------

% در ورژن جدید زی‌پرشین برای تایپ متن‌های ریاضی، این سه بسته، حتماً باید فراخوانی شود
\usepackage{amsthm,amssymb,amsmath}
% بسته‌ای برای تنطیم حاشیه‌های بالا، پایین، چپ و راست صفحه
\usepackage[top=30mm, bottom=25mm, left=25mm, right=35mm]{geometry}
% بسته‌‌ای برای ظاهر شدن شکل‌ها و تصاویر متن
\usepackage{graphicx}
% بسته‌ای برای رسم کادر
\usepackage{framed} 
% بسته‌‌ای برای چاپ شدن خودکار تعداد صفحات در صفحه «معرفی پایان‌نامه»
\usepackage{lastpage}
% بسته‌ و دستوراتی برای ایجاد لینک‌های رنگی با امکان جهش
% \usepackage[pagebackref=false,colorlinks,linkcolor=black,citecolor=black]{hyperref}
% چنانچه قصد پرینت گرفتن نوشته خود را دارید، خط بالا را غیرفعال و  از دستور زیر استفاده کنید چون در صورت استفاده از دستور زیر‌‌، 
% لینک‌ها به رنگ سیاه ظاهر خواهند شد که برای پرینت گرفتن، مناسب‌تر است
% \usepackage[pagebackref=false]{hyperref}
% بسته‌ لازم برای تنظیم سربرگ‌ها
\usepackage{fancyhdr}
\setlength{\headheight}{16pt}
%
\usepackage{setspace}
\usepackage{algorithm}
\usepackage{algorithmic}
\usepackage{subfigure}
\usepackage{titletoc}


% بسته‌ای برای ظاهر شدن «مراجع» و «نمایه» در فهرست مطالب
\usepackage[nottoc]{tocbibind}
% دستورات مربوط به ایجاد نمایه
\usepackage{makeidx}
\makeindex
% تعریف شیم‌های مورد نیاز برای نسخه‌های جدید بسته bidi در برخی محیط‌ها
% Removal of redundant and potentially broken definitions
\makeatletter
% شیم‌های سازگاری برای نسخه‌های جدید لاتک و بسته‌های bidi/xepersian
\providecommand{\IfClassLoadedT}[2]{\@ifclassloaded{#1}{#2}{}}
\providecommand{\IfClassLoadedF}[2]{\@ifclassloaded{#1}{}{#2}}
\providecommand{\IfClassLoadedTF}[3]{\@ifclassloaded{#1}{#2}{#3}}
\providecommand{\IfPackageLoadedT}[2]{\@ifpackageloaded{#1}{#2}{}}
\providecommand{\IfPackageLoadedF}[2]{\@ifpackageloaded{#1}{}{#2}}
\providecommand{\IfPackageLoadedTF}[3]{\@ifpackageloaded{#1}{#2}{#3}}
\providecommand{\@iffileloaded}[3]{\expandafter\ifx\csname ver@#1\endcsname\relax#3\else#2\fi}
\providecommand{\IfFileLoadedT}[2]{\@iffileloaded{#1}{#2}{}}
\providecommand{\IfFileLoadedF}[2]{\@iffileloaded{#1}{}{#2}}
\providecommand{\IfFileLoadedTF}[3]{\@iffileloaded{#1}{#2}{#3}}
\makeatother

%%%%%%%%%%%%%%%%%%%%%%%%%%
% فراخوانی بسته زی‌پرشین و تعریف قلم فارسی و انگلیسی
\usepackage[pagebackref=false,colorlinks,linkcolor=black,citecolor=black]{hyperref}
\usepackage{notoccite}

\ExplSyntaxOn
\cs_if_exist:NTF \g__hyp_linknestlevel_int { } { \int_new:N \g__hyp_linknestlevel_int }
\cs_if_exist:NTF \l__hyp_annot_GoTo_bool { } { \bool_new:N \l__hyp_annot_GoTo_bool }
\cs_if_exist:NTF \l__hyp_annot_URI_bool { } { \bool_new:N \l__hyp_annot_URI_bool }
\cs_if_exist:NTF \l__hyp_href_url_ismap_bool { } { \bool_new:N \l__hyp_href_url_ismap_bool }
\cs_if_exist:NTF \l__hyp_text_enc_uri_print_tl { } { \tl_new:N \l__hyp_text_enc_uri_print_tl }
\cs_if_exist:NTF \l__hyp_uri_tmpa_tl { } { \tl_new:N \l__hyp_uri_tmpa_tl }

% Defining missing l3pdf/pdfdict commands if not present
\cs_if_exist:NTF \pdfdict_put:nno { } { \cs_new:Npn \pdfdict_put:nno #1#2#3 { } }
\cs_if_exist:NTF \pdfdict_put:nnn { } { \cs_new:Npn \pdfdict_put:nnn #1#2#3 { } }
\cs_if_exist:NTF \pdfdict_use:n { } { \cs_new:Npn \pdfdict_use:n #1 { } }
\cs_if_exist:NTF \pdfannot_dict_put:nne { } { \cs_new:Npn \pdfannot_dict_put:nne #1#2#3 { } }
\cs_if_exist:NTF \pdfannot_link:nen { } { \cs_new:Npn \pdfannot_link:nen #1#2#3 { } }
\cs_if_exist:NTF \__hyp_check_link_nesting:TF { } { \cs_new:Npn \__hyp_check_link_nesting:TF #1#2 { #1 } }
\cs_if_exist:NTF \__hyp_text_pdfstring:eoN { } { \cs_new:Npn \__hyp_text_pdfstring:eoN #1#2#3 { } }
\ExplSyntaxOff

% تعریف دستورات گمشده برای جلوگیری از خطای renewcommand در بسته‌های داخلی
\makeatletter
\providecommand{\foottextfont}{}
\providecommand{\LTRfoottextfont}{}
\providecommand{\RTLfoottextfont}{}
\providecommand{\abstractname}{}
\providecommand{\latinabstract}{}
\providecommand{\AutoPersianDigits}{}
\makeatother

\usepackage{xepersian}
\settextfont[Path=font/,Extension=.ttf,Scale=1,
    BoldFont=XB NiloofarBd,
    ItalicFont=XB NiloofarIt,
    BoldItalicFont=XB NiloofarBdIt
]{XB Niloofar}
\setlatintextfont[Scale=0.9]{Times New Roman}

%%%%%%%%%%%%%%%%%%%%%%%%%%
% چنانچه می‌خواهید اعداد در فرمول‌ها، انگلیسی باشد، خط زیر را غیرفعال کنید
\setdigitfont[Path=font/,Extension=.ttf,Scale=1,
    BoldFont=XB NiloofarBd,
    ItalicFont=XB NiloofarIt,
    BoldItalicFont=XB NiloofarBdIt
]{XB Niloofar}
%%%%%%%%%%%%%%%%%%%%%%%%%%
% تعریف قلم‌های فارسی و انگلیسی اضافی برای استفاده در بعضی از قسمت‌های متن
\defpersianfont\titlefont[Path=font/,Extension=.ttf,Scale=1]{XB Titre}
% \defpersianfont\iranic[Scale=1.1]{XB Zar Oblique}%Italic}%
% \defpersianfont\nastaliq[Scale=1.2]{IranNastaliq}

%%%%%%%%%%%%%%%%%%%%%%%%%%
% دستوری برای حذف کلمه «چکیده»
\renewcommand{\abstractname}{}
% دستوری برای حذف کلمه «abstract»
%\renewcommand{\latinabstract}{}
% دستوری برای تغییر نام کلمه «اثبات» به «برهان»
\renewcommand\proofname{\textbf{برهان}}
% دستوری برای تغییر نام کلمه «کتاب‌نامه» به «مراجع»
\renewcommand{\bibname}{مراجع}
% دستوری برای تعریف واژه‌نامه انگلیسی به فارسی
\newcommand\persiangloss[2]{#1\dotfill\lr{#2}\\}
% دستوری برای تعریف واژه‌نامه فارسی به انگلیسی 
\newcommand\englishgloss[2]{#2\dotfill\lr{#1}\\}
% تعریف دستور جدید «\پ» برای خلاصه‌نویسی جهت نوشتن عبارت «پروژه/پایان‌نامه/رساله»
\newcommand{\پ}{پروژه/پایان‌نامه/رساله }

%\newcommand\BackSlash{\char`\\}

%%%%%%%%%%%%%%%%%%%%%%%%%%
\SepMark{-}

% تعریف و نحوه ظاهر شدن عنوان قضیه‌ها، تعریف‌ها، مثال‌ها و ...
\theoremstyle{definition}
\newtheorem{definition}{تعریف}[section]
\theoremstyle{plain}
\newtheorem{theorem}[definition]{قضیه}
\newtheorem{lemma}[definition]{لم}
\newtheorem{proposition}[definition]{گزاره}
\newtheorem{corollary}[definition]{نتیجه}
\newtheorem{remark}[definition]{ملاحظه}
\theoremstyle{definition}
\newtheorem{example}[definition]{مثال}

%\renewcommand{\theequation}{\thechapter-\arabic{equation}}
%\def\bibname{مراجع}
\numberwithin{algorithm}{chapter}
\def\listalgorithmname{فهرست الگوریتم‌ها}
\def\listfigurename{فهرست تصاویر}
\def\listtablename{فهرست جداول}

%%%%%%%%%%%%%%%%%%%%%%%%%%%%
% دستورهایی برای سفارشی کردن سربرگ صفحات
% \newcommand{\SetHeader}{
% \csname@twosidetrue\endcsname
% \pagestyle{fancy}
% \fancyhf{} 
% \fancyhead[OL,EL]{\thepage}
% \fancyhead[OR]{\small\rightmark}
% \fancyhead[ER]{\small\leftmark}
% \renewcommand{\chaptermark}[1]{%
% \markboth{\thechapter-\ #1}{}}
% }
%%%%%%%%%%%%5
%\def\MATtextbaseline{1.5}
%\renewcommand{\baselinestretch}{\MATtextbaseline}
\doublespacing
%%%%%%%%%%%%%%%%%%%%%%%%%%%%%
% دستوراتی برای اضافه کردن کلمه «فصل» در فهرست مطالب با استفاده از titletoc
\titlecontents{chapter}
  [5em]
  {\addvspace{1em}\bfseries}
  {\contentslabel[\chaptername~\thecontentslabel:]{5em}}
  {\hspace*{-5em}}
  {\titlerule*[0.5pc]{.}\contentspage}

\titlecontents{section}
  [5em]
  {}
  {\contentslabel{2em}}
  {\hspace*{-2em}}
  {\titlerule*[0.5pc]{.}\contentspage}

\titlecontents{subsection}
  [8em]
  {}
  {\contentslabel{3em}}
  {\hspace*{-3em}}
  {\titlerule*[0.5pc]{.}\contentspage}

\titlecontents{subsubsection}
  [12em]
  {}
  {\contentslabel{4em}}
  {\hspace*{-4em}}
  {\titlerule*[0.5pc]{.}\contentspage}

\makeatletter
{\def\thebibliography#1{\chapter*{\refname\@mkboth
   {\uppercase{\refname}}{\uppercase{\refname}}}\list
   {[\arabic{enumi}]}{\settowidth\labelwidth{[#1]}
   \rightmargin\labelwidth
   \advance\rightmargin\labelsep
   \advance\rightmargin\bibindent
   \itemindent -\bibindent
   \listparindent \itemindent
   \parsep \z@
   \usecounter{enumi}}
   \def\newblock{}
   \sloppy
   \sfcode`\.=1000\relax}}
\makeatother


% تنظیمات اختصاصی برای فهرست مطالب (TOC)
\setcounter{secnumdepth}{3}
\setcounter{tocdepth}{3}

% بازگرداندن نقاط رابط (Dot Leaders) که در Cleanup حذف شده بود
\renewcommand{\cftchapleader}{\cftdotfill{\cftdotsep}}
\renewcommand{\cftsecleader}{\cftdotfill{\cftdotsep}}
\renewcommand{\cftsubsecleader}{\cftdotfill{\cftdotsep}}

% تنظیمات تورفتگی و ظاهر
\setlength{\cftsubsubsecindent}{6em}
\setlength{\cftaftertoctitleskip}{10pt} % رفع فاصله زیاد بعد از عنوان فهرست
\raggedbottom

\begin{document}
\AutoPersianDigits
\setstretch{1.5}

\pagenumbering{harfi}
\firstPage
\besmPage
\davaranPage
% در این قسمت اسامی‌داوران در فایل اصلی باید دستی وارد شوند، اینجا از ساختار پیش‌فرض استفاده می‌کنیم

\esalatPage
\mojavezPage

\chapter*{تقدیم}
\addcontentsline{toc}{chapter}{تقدیم}
\begin{center}
    \textit{(متن تقدیم به خانواده، اساتید و ...)}
\end{center}
\newpage

\begin{acknowledgementpage}
(متن تشکر و قدردانی از اساتید راهنما، مشاور و دوستان)
\signature 
\end{acknowledgementpage}

\keywords{سیستم‌های خودوفقی، یادگیری تقویتی عمیق، شکل‌دهی پاداش، عدم قطعیت}
\fa-abstract{
(خلاصه‌ای از چالش، راهکار و نتایج کلیدی در ۳۰۰ تا ۵۰۰ کلمه)
}
\abstractPage

\tableofcontents
\newpage

\listoftables
\newpage

\listoffigures
\newpage

\chapter*{فهرست علائم و اختصارات}
\addcontentsline{toc}{chapter}{فهرست علائم و اختصارات}
(لیست نمادهای ریاضی مانند $\rho$، $\gamma$ و اختصارات مانند \lr{DRL}، \lr{SAS})
\newpage

% ============================================================================
% متن اصلی رساله
% ============================================================================
\pagenumbering{arabic}
\pagestyle{fancy}

% فصل ۱: مقدمه
\chapter{مقدمه}
\section{پیش‌زمینه و مفاهیم: سیستم‌های خودوفقی}
\subsection{انگیزه‌های پژوهش و مثال‌های عینی}
\subsection{حلقه بازخورد \lr{MAPE-K}}
\subsection{کاستی‌های رویکردهای ایستا و مبتنی بر قانون}

\section{چالش عدم قطعیت در زمان اجرا}
\subsection{مبانی عدم قطعیت در سیستم‌های خودوفقی}
\subsection{مطالعه موردی انگیزشی: سیستم \lr{DeltaIoT} (نگاه اول)}
\subsection{منابع و اثرات عدم قطعیت در \lr{DeltaIoT}}

\section{یادگیری تقویتی برای خودوفقی: فرصت‌ها و چالش‌ها}
\subsection{چرا یادگیری تقویتی برای تصمیم‌گیری در زمان اجرا؟}
\subsection{محدودیت‌های یادگیری تقویتی سنتی در فضاهای وفقی گسسته و بزرگ}

\section{بیان مسئله و سوالات تحقیق}

\section{راهکار پیشنهادی: یادگیری تقویتی عمیق با شکل‌دهی پاداش (\lr{RS-DRL})}
\textit{ادغام شده در معماری حلقه بازخورد \lr{MAPE-K} (\lr{Monitor-Analyze-Plan-Execute-Knowledge})}

\subsection{ایده اصلی: شکل‌دهی تصادفی پاداش (فاکتور \texorpdfstring{$\rho$}{rho}) در طول بازپخش تجربه}
\textit{کنترل شده توسط یک فاکتور شکل‌دهی \texorpdfstring{$\rho$}{rho} که نسبت تجربیات ناموفق برای شکل‌دهی مجدد را تعیین می‌کند.}

\section{مشارکت‌های کلیدی رساله}
\section{ساختار کلی رساله}

\newpage

% فصل ۲: مبانی و پیشینه تحقیق
\chapter{مبانی و پیشینه تحقیق}
\section{مبانی سیستم‌های خودوفقی}
\subsection{مفاهیم اصلی و الگوهای معماری}
\subsection{مدل‌سازی و مدیریت عدم قطعیت در سیستم‌های خودوفقی}

\section{یادگیری تقویتی برای وفقی‌سازی سیستم}
\subsection{اصول یادگیری تقویتی (\lr{MDP}ها، یادگیری \lr{Q})}
\subsection{یادگیری تقویتی عمیق برای فضاهای عمل گسسته بزرگ (\lr{DQN})}
\subsection{مفهوم یادگیری «مداوم» در این رساله (حلقه‌های وفقی جریانی)}

\section{کار‌های مرتبط: شیوه‌های مدیریت فضاهای وفقی بزرگ}
\subsection{روش‌های مبتنی بر یادگیری ماشین (به‌عنوان مثال \lr{ML4EAS})}
\subsection{روش‌های مبتنی بر جستجو (به‌عنوان مثال الگوریتم‌های ژنتیک)}
\subsection{روش‌های مبتنی بر یادگیری تقویتی (مانند \lr{Epsilon-Greedy DQN} و \texorpdfstring{\\}{}\lr{DLASER})}
\subsection{شکل‌دهی پاداش و بازپخش تجربه با نگاه به گذشته (\lr{HER} و \lr{Soft HER})}

\section{سنتز و شناسایی شکاف تحقیقاتی}

\newpage

% فصل ۳: مدل سیستم و فرمول‌بندی مسئله
\chapter{مدل سیستم و فرمول‌بندی مسئله}
\section{مطالعه موردی انگیزشی: سیستم \lr{DeltaIoT}}
\subsection{مرور کلی سیستم و منابع عدم قطعیت}
\subsection{فضای وفقی و بازنمایی پیکربندی}

\section{مدل صوری برای تصمیم‌گیری خودوفقی تحت عدم قطعیت}
\subsection{تعریف فضای حالت}
\textit{فضای حالت به صورت $S = \{\text{\lr{Energy}}, \text{\lr{Packet Loss}}, \text{\lr{Latency}}\}$ تعریف می‌شود.}
\subsection{فضای عمل: گزینه‌های وفقی گسسته}
\subsection{فرمول‌بندی تابع پاداش برای اهداف تک‌هدفه و چندهدفه}
\subsubsection{معادله ۱: پاداش چندهدفه (معادل معادله ۱۱ مقاله)}
\subsubsection{معادله ۲: پاداش تک‌هدفه (معادل معادله ۱۲ مقاله)}
\subsubsection{جدول ۱: پارامترهای پاداش برای اهداف \lr{TTS} و \lr{TT} (معادل جدول ۲ مقاله)}

\section{بیان مسئله: یادگیری سیاست وفقی بهینه با عملیات مداوم در زمان اجرا}

\newpage

% فصل ۴: روش \lr{RS-DRL}: معماری و الگوریتم
\chapter{روش \lr{RS-DRL}: معماری و الگوریتم}
\section{مرور کلی رویکرد پیشنهادی}

\section{معماری سیستم: ادغام \lr{RS-DRL} با حلقه \lr{MAPE-K}}
\subsection{سیستم مدیریت‌شده و سیستم مدیریت‌کننده}
\subsection{تعامل دقیق حلقه وفقی ۹ مرحله‌ای}

\section{الگوریتم \lr{RS-DRL}}
\subsection{الگوریتم اصلی: \lr{DQN} برای حلقه‌های وفقی جریانی}
\textit{رفع محدودیت‌های یادگیری تقویتی کلاسیک از طریق تعامل یادگیری جریانی و مداوم.}
\subsection{تابع شکل‌دهی پاداش (الگوریتم ۲)}
\subsection{مقایسه با بازپخش تجربه با نگاه به گذشته (\lr{HER}) و سایر روش‌های پایه}

\section{جزئیات کلیدی پیاده‌سازی برای بازتولیدپذیری}
\subsection{ساخت حالت از داده‌های پایش شده}
\subsection{اعتبارسنجی زمان اجرا با \lr{UPPAAL-SMC} و \lr{ActivFORMS}}
\subsection{انتخاب و بهینه‌سازی ابرپارامترها (جستجوی شبکه‌ای و تنظیم \texorpdfstring{$\rho$}{rho})}

\newpage

% فصل ۵: ارزیابی تجربی
\chapter{ارزیابی تجربی}
\section{اهداف ارزیابی و سوالات تحقیق}

\section{تنظیمات آزمایشگاهی}
\subsection{نمونه‌های مطالعه موردی: \lr{DeltaIoTv1}، \lr{DeltaIoTv1.1} و \lr{DeltaIoTv2}}
\subsection{مجموعه داده‌ها: ردپاهای محیطی و پارامترهای رانش}
\subsection{پلتفرم پیاده‌سازی: \lr{Gymnasium} و \lr{Stable-Baselines3}}

\section{روش‌های پایه و مقایسه}
\subsection{روش \lr{Epsilon-Greedy DQN}}
\subsection{روش \lr{RS-DRL} به همراه \lr{HER / Soft HER}}
\subsection{روش‌های \lr{DLASER}، \lr{DLASER+} و \lr{ML4EAS}}
\subsection{روش مرجع جستجوی کامل (\lr{Exhaustive Search})}
\subsection{توجیه انتخاب روش‌های پایه: محدودیت‌ها و پروتکل‌های مقایسه}

\section{معیارهای ارزیابی}
\subsection{عملکرد یادگیری: عملکرد مجانبی، زمان تا آستانه و عملکرد کل}
\subsection{عملکرد سیستم: از دست رفتن بسته، تاخیر، مصرف انرژی}

\section{نتایج و تحلیل‌ها}
\subsection{منحنی‌های یادگیری و رفتار همگرایی}
\subsection{مقایسه عملکرد به ازای هر هدف}
\subsection{ارزیابی مقایسه‌ای در برابر تمام روش‌های پایه}
\subsection{مطالعه ابلیشن (\lr{Ablation}): حساسیت به فاکتور شکل‌دهی (\texorpdfstring{$\rho$}{rho})}
\subsection{تحلیل مقیاس‌پذیری (از ۲۱۶ به ۴۰۹۶ گزینه وفقی)}

\section{اعتبارسنجی آماری (آزمون‌های \lr{t} و مقادیر \lr{p})}

\newpage

% فصل ۶: بحث، کارهای آینده و نتیجه‌گیری
\chapter{بحث، کارهای آینده و نتیجه‌گیری}
\section{بحث در مورد یافته‌های کلیدی}
\subsection{نظریه‌پردازی مکانیزم: چرا \lr{RS-DRL} کار می‌کند؟}
\subsection{تحلیل تعاملات و مصالحه‌ها (به‌عنوان مثال، تاخیر در مقابل از دست رفتن بسته)}
\subsection{پیامدهای مهندسی سیستم‌های خودوفقی}

\section{محدودیت‌ها و تهدیدات اعتبار}
\subsection{اعتبار داخلی، خارجی و سازه}
\subsection{محدودیت‌های رویکرد \lr{RS-DRL} (وابستگی به \texorpdfstring{$\rho$}{rho}، فضاهای عمل گسسته)}

\section{مسیرهایی برای تحقیقات آینده}
\subsection{به سوی استراتژی‌های شکل‌دهی وفقی (تنظیم خودکار \texorpdfstring{$\rho$}{rho})}
\subsection{گسترش \lr{RS-DRL} به دامنه‌های پیوسته و حساس به ایمنی}
\subsection{اعتبارسنجی بین‌دامنه و تحلیل‌های نظری}

\section{نتیجه‌گیری}
\subsection{خلاصه بیان مسئله و راهکار پیشنهادی}
\subsection{مرور مشارکت‌ها و شواهد تجربی}
\subsection{سخنان پایانی}

\newpage

% کتابنامه و پیوست‌ها
\pagestyle{empty}
{
\onehalfspacing
\bibliographystyle{ieeetr-fa}
\nocite{*}
\bibliography{MyReferences}
}

\pagestyle{fancy}
\appendix
\chapter{جریان‌های کاری کامل فرآیند}
\newpage
\chapter{لیست کامل الگوریتم‌ها}
\newpage
\chapter{نتایج تجربی گسترده}
\newpage
\chapter{مصنوعات پیاده‌سازی و پکیج بازتولید}

\newpage

% چکیده و عنوان انگلیسی
\en-abstract{
(English summary of the thesis, required at the end by TMU guidelines)
}
\latinfirstPage

\end{document}
